% Risposta alla sezione 1.2 di Exercise01.pdf.
\documentclass[11pt,a4paper]{article}
\usepackage[T1]{fontenc}
\usepackage[utf8]{inputenc}
\usepackage[english]{babel}
\usepackage{lmodern}
\usepackage[margin=2.2cm]{geometry}
\usepackage{amsmath,amssymb}
\usepackage{booktabs}
\usepackage{microtype}
\usepackage[hidelinks]{hyperref}
\setlength{\parindent}{0pt}
\setlength{\parskip}{0.55em}
\setlength{\emergencystretch}{2em}
\title{Exercise 1.2: Amdahl's Law}
\author{}
\date{}
\begin{document}
\maketitle
\vspace{-2em}
\section*{1. Derivation}
Let $T_1$ be the sequential execution time of a fixed workload, $P$ the fraction
of this time that can be parallelized, and $N$ the number of processors.
The serial part takes $(1-P)T_1$. Assuming perfect load balancing, no parallel
overhead, and an ideal $N$-fold acceleration of the parallel part, its execution
time is $PT_1/N$. Hence
\begin{equation}
  T_N=(1-P)T_1+\frac{PT_1}{N}
     =T_1\left((1-P)+\frac{P}{N}\right).
\end{equation}
The overall speedup is therefore
\begin{equation}
  \boxed{S(N)=\frac{T_1}{T_N}=\frac{1}{(1-P)+P/N}}.
  \label{eq:law}
\end{equation}
For $0\leq P<1$,
\begin{equation}
  \lim_{N\to\infty}S(N)=\frac{1}{1-P}.
\end{equation}
The serial fraction sets the asymptotic limit. For the special case $P=1$,
the ideal speedup is $S(N)=N$ and there is no finite serial-fraction limit.

\section*{2. Correctness for real applications}
Amdahl's law is correct under its assumptions. It is not a universal prediction
for every parallel execution: using $N$ as the acceleration of the parallel
portion assumes that overhead and other architectural effects can be neglected.

\paragraph{Performance predicted too high.}
Communication, synchronization, thread creation, and load imbalance increase
the actual time. Consider a computation with an ideal $100\times$ speedup
whose processors exchange intermediate results
over a network. If communication becomes a bottleneck, the actual speedup
is below $100\times$. Writing
\[
  T_N=T_1\left((1-P)+\frac{P}{N}\right)+T_{\mathrm{overhead}}(N)
\]
makes clear why positive overhead reduces the ideal speedup.

\paragraph{Performance predicted too low.}
A working set that exceeds the cache of one processor may fit into the aggregate caches
when partitioned across several processors. Reducing slow memory accesses
can accelerate the parallel part by more than $N$, so actual speedup can
exceed the prediction of Equation~\eqref{eq:law}, and can even be superlinear.
This changes the assumed acceleration of the parallel portion; it does not
invalidate the time-decomposition argument.

\paragraph{Changing the workload.}
In weather forecasting, more processors are often used to increase grid
resolution rather than to finish an unchanged
problem sooner. Amdahl's fixed-workload limit does not directly characterize
this comparison. Weak scaling studies increasing the workload in proportion
to the number of processors, while Gustafson's law addresses scaled speedup
using the serial fraction of a parallel run. These perspectives change the
workload assumptions used in the fixed-size Amdahl model.

\section*{3. Numerical calculations and comparison}
\subsection*{(a) $P=0.5$ and $N=4$}
\[
  S_A=\frac{1}{(1-0.5)+0.5/4}
     =\frac{1}{0.5+0.125}
     =\boxed{1.6}.
\]
The parallel execution takes $0.625T_1$, and its efficiency is
$S_A/4=0.4=40\%$.

\subsection*{(b) $P=0.2$ and $N=8$}
\[
  S_B=\frac{1}{(1-0.2)+0.2/8}
     =\frac{1}{0.8+0.025}
     =\boxed{\frac{40}{33}\approx1.212}.
\]
The parallel execution takes $0.825T_1$, and its efficiency is
$S_B/8=5/33\approx15.15\%$.

\subsection*{(c) Comparison}
\begin{center}
\begin{tabular}{@{}lrrrr@{}}
\toprule
Case & $P$ & $N$ & Speedup & Limit as $N\to\infty$\\
\midrule
A & $0.5$ & $4$ & $1.6$ & $2$\\
B & $0.2$ & $8$ & $1.212\ldots$ & $1.25$\\
\bottomrule
\end{tabular}
\end{center}
Case A is faster relative to its sequential baseline despite having fewer
processors. Its larger parallelizable fraction leaves a smaller serial
bottleneck. In case B, $80\%$ of the original time remains serial, so even
arbitrarily many processors cannot produce a speedup above $1.25$.

These cases change both $P$ and $N$, so their difference cannot be attributed
to processor count alone. For any fixed $P>0$, adding processors increases
ideal speedup, but with diminishing returns when $P<1$.
Reducing the serial fraction is essential for effective scaling.

For comparison, keeping $P=0.5$ and increasing only $N$ from four to eight
would give
\[
  S(8)=\frac{1}{0.5+0.5/8}=\frac{16}{9}\approx1.778.
\]

\end{document}
